\documentclass[a4paper,12pt]{article}
\usepackage[T2A]{fontenc}
\usepackage[utf8]{inputenc}
\usepackage[russian]{babel}
\usepackage{amsmath,amssymb,mathtools}
\usepackage{helvet}
\renewcommand{\familydefault}{\sfdefault}
\usepackage{icomma}
\usepackage[a4paper,margin=18mm]{geometry}
\usepackage{microtype}
\usepackage{tikz}
\usetikzlibrary{arrows.meta,positioning,shapes.geometric,calc}
\usepackage{tabularx,booktabs,longtable}
\usepackage{enumitem}
\usepackage{hyperref}
\usepackage{parskip}
\usepackage{fancyhdr}
\usepackage{setspace}
\usepackage{xcolor}

\definecolor{accent}{HTML}{285F73}
\definecolor{darkaccent}{HTML}{173C49}
\definecolor{textgray}{HTML}{333333}
\definecolor{softaccent}{HTML}{DCE9ED}
\definecolor{warm}{HTML}{A85D4A}

\hypersetup{colorlinks=true,linkcolor=darkaccent,urlcolor=darkaccent}
\setlength{\parindent}{0pt}
\setlength{\headheight}{15pt}
\setlength{\parskip}{0.55em}
\onehalfspacing
\pagestyle{fancy}
\fancyhf{}
\lhead{\small Текстовые задачи на движение}
\rhead{\small 28.09.26}
\cfoot{\thepage}
\renewcommand{\headrulewidth}{0.3pt}
\setlist[itemize]{leftmargin=1.4em,itemsep=0.25em,topsep=0.3em}
\setlist[enumerate]{leftmargin=1.6em,itemsep=0.45em,topsep=0.4em}

\tikzset{
  flowstart/.style={ellipse,draw=accent,line width=.9pt,align=center,minimum width=34mm,minimum height=10mm,font=\small},
  flowaction/.style={rectangle,rounded corners=3pt,draw=accent,line width=.9pt,align=center,text width=54mm,minimum height=12mm,font=\small},
  flowchoice/.style={diamond,aspect=2.1,draw=warm,line width=.9pt,align=center,text width=28mm,inner sep=1.6pt,font=\small},
  flowbranch/.style={rectangle,rounded corners=3pt,draw=accent,line width=.9pt,align=center,text width=41mm,minimum height=12mm,font=\small},
  flowarrow/.style={-{Latex[length=2.5mm]},line width=.8pt,color=textgray}
}

\newcommand{\vts}{\(v\)--\(t\)--\(S\)}
\newcommand{\lessonrule}{\vspace{0.25em}\hrule\vspace{0.65em}}

\begin{document}

\thispagestyle{empty}
\begin{center}
\vspace*{18mm}
{\color{accent}\Large Чек-лист}\par
\vspace{5mm}
{\Huge\bfseries Текстовые задачи на движение}\par
\vspace{3mm}
{\Large круговая трасса · движение по воде · средняя скорость}\par
\vspace{12mm}
{\large Анастасия Павлова}\par
\vspace{3mm}
{\normalsize 28 сентября 2026}\par
\vfill
{\small Лёвин Артём Александрович эксклюзивно для Анастасии Павловой}\par
\vspace{9mm}
\begin{tikzpicture}[x=1cm,y=1cm]
  \draw[accent,line width=1.2pt]
    (-7,0) .. controls (-4.8,0.9) and (-2.3,-0.8) .. (0,0)
           .. controls (2.2,0.8) and (4.6,-0.7) .. (7,0.1);
  \fill[warm] (-4.7,0.55) circle (1.4pt);
  \fill[accent] (0,0) circle (1.4pt);
  \fill[warm] (4.7,-0.35) circle (1.4pt);
\end{tikzpicture}
\end{center}
\clearpage

\section*{1. Универсальная опора: таблица \vts}

Для равномерного движения используются три величины:
\[
S=vt,\qquad v=\frac{S}{t},\qquad t=\frac{S}{v}.
\]

\begin{tabularx}{\linewidth}{@{}lXX@{}}
\toprule
\textbf{Величина} & \textbf{Что означает} & \textbf{Главный вопрос при заполнении таблицы} \\
\midrule
\(v\) & скорость & известна ли она напрямую или выражается через неизвестную? \\
\(t\) & время & одинаково ли оно у участников, отличается ли на известную величину? \\
\(S\) & путь & одинаков ли путь, складываются ли пути или берётся их разность? \\
\bottomrule
\end{tabularx}

\textbf{Рабочая последовательность:}
\begin{enumerate}
  \item Определите, что требуется найти, и удобно обозначьте это буквой.
  \item Разделите ситуацию на строки таблицы: участники, участки пути, движение по/против течения, стоянка.
  \item Заполните известные величины и один столбец вычислите по формулам движения.
  \item Составьте \textbf{смысловое} уравнение из условия: сумма, разность, одинаковое время, полный путь и т.~п.
  \item Решите уравнение и проверьте ответ по смыслу и единицам измерения.
\end{enumerate}

\lessonrule
\textbf{Единицы измерения.} Перед составлением уравнения приведите время и расстояние к совместимым единицам. Например,
\[
40\text{ мин}=\frac{2}{3}\text{ ч}.
\]

\section*{2. Движение по круговой трассе}

Если два участника движутся \textbf{в одном направлении}, быстрый постепенно набирает путь относительно медленного. Поэтому основная величина — \textbf{разность путей}:
\[
(v_{\text{быстр}}-v_{\text{медл}})t=\Delta S.
\]

\textbf{Старт из одной точки, обгон на один круг:}
\[
\Delta S=L,
\]
где \(L\) — длина круга.

\textbf{Старт из диаметрально противоположных точек:}
если быстрый должен догнать медленного и начальный разрыв равен половине круга, то
\[
\Delta S=\frac{L}{2}.
\]

\subsection*{Пример 1. Обгон на один круг}
Длина круга \(15\) км, скорости \(80\) и \(60\) км/ч. Тогда
\[
80t-60t=15,
\qquad 20t=15,
\qquad t=\frac34\text{ ч}=45\text{ мин}.
\]

\subsection*{Пример 2. Половина окружности}
Длина круга \(20\) км, разность скоростей равна \(12\) км/ч, участники стартуют из диаметрально противоположных точек, быстрый находится сзади по ходу движения. Тогда
\[
12t=\frac{20}{2}=10,
\qquad t=\frac56\text{ ч}=50\text{ мин}.
\]

\textbf{Типичная ошибка:} складывать пути при движении в одном направлении. Для догоняющего движения сравнивается именно \textbf{разность} пройденных путей.

\section*{3. Движение по воде}

Пусть \(v\) — собственная скорость лодки в неподвижной воде, \(u\) — скорость течения. Тогда
\[
v_{\text{по}}=v+u,
\qquad
v_{\text{против}}=v-u.
\]
Обязательное условие физического смысла:
\[
v>u>0.
\]

Для участка длины \(S\):
\[
t_{\text{по}}=\frac{S}{v+u},
\qquad
t_{\text{против}}=\frac{S}{v-u}.
\]

\subsection*{Пример 3. Разность времён}
Путь в каждую сторону — \(24\) км, скорость течения — \(2\) км/ч. На путь против течения требуется на \(3\) ч больше. Ищем собственную скорость \(v\):
\[
\frac{24}{v-2}-\frac{24}{v+2}=3,
\qquad v>2.
\]
Делим уравнение на \(3\):
\[
\frac{8}{v-2}-\frac{8}{v+2}=1.
\]
После приведения к общему знаменателю:
\[
\frac{32}{v^2-4}=1,
\qquad v^2=36.
\]
С учётом \(v>2\):
\[
\boxed{v=6\text{ км/ч}}.
\]

\subsection*{Стоянка и полный рейс}
Стоянка — отдельная строка таблицы. Если путь туда равен \(x\), обратно тоже \(x\), то полный путь равен \(2x\).

Например, скорости по и против течения равны \(24\) и \(16\) км/ч, стоянка длится \(4\) ч, весь рейс — \(14\) ч:
\[
\frac{x}{24}+4+\frac{x}{16}=14.
\]
Отсюда \(x=96\) км, поэтому полный путь:
\[
\boxed{2x=192\text{ км}}.
\]

\subsection*{Плот}
У плота нет собственной скорости относительно воды, поэтому его скорость относительно берега равна скорости течения:
\[
v_{\text{плота}}=u.
\]
Плот движется по течению; в стандартной школьной модели движение плота против течения не рассматривается.

\textbf{Типичные ошибки:}
\begin{itemize}
  \item перепутать знаки: по течению \(v+u\), против течения \(v-u\);
  \item забыть ограничение \(v>u\);
  \item не выделить стоянку отдельным временем;
  \item найти путь в одну сторону, когда в вопросе требуется весь рейс.
\end{itemize}

\clearpage
\section*{Алгоритм решения текстовой задачи на движение}
\thispagestyle{plain}
\vspace{4mm}
\begin{center}
\begin{tikzpicture}[node distance=8mm and 12mm]
  \node[flowstart] (start) {Прочитать условие\\и определить искомое};
  \node[flowaction,below=of start] (table) {Разбить движение на случаи и заполнить таблицу \(v\)--\(t\)--\(S\)};
  \node[flowchoice,below=10mm of table] (type) {Какой смысловой тип?};

  \node[flowbranch,below left=14mm and 26mm of type] (circle) {Круг/догонка:\\разность путей = нужная дуга};
  \node[flowbranch,below=14mm of type] (water) {Вода:\\\(v+u\), \(v-u\); стоянку учитывать отдельно};
  \node[flowbranch,below right=14mm and 26mm of type] (avg) {Средняя скорость:\\весь путь / всё время};

  \node[flowaction,below=27mm of water] (eq) {Составить смысловое уравнение и указать ограничения};
  \node[flowaction,below=of eq] (solve) {Решить уравнение, проверить единицы и смысл ответа};
  \node[flowstart,below=of solve] (finish) {Записать ответ};

  \draw[flowarrow] (start) -- (table);
  \draw[flowarrow] (table) -- (type);
  \draw[flowarrow] (type) -- node[above left,font=\scriptsize]{круг} (circle);
  \draw[flowarrow] (type) -- node[right,font=\scriptsize]{вода} (water);
  \draw[flowarrow] (type) -- node[above right,font=\scriptsize]{средняя} (avg);
  \draw[flowarrow] (circle) |- (eq);
  \draw[flowarrow] (water) -- (eq);
  \draw[flowarrow] (avg) |- (eq);
  \draw[flowarrow] (eq) -- (solve);
  \draw[flowarrow] (solve) -- (finish);
\end{tikzpicture}
\end{center}
\clearpage

\section*{4. Средняя скорость}

Средняя скорость на всём маршруте определяется только через \textbf{полный путь} и \textbf{полное время}:
\[
\boxed{v_{\text{ср}}=\frac{S_{\text{общ}}}{t_{\text{общ}}}}
=\frac{S_1+S_2+\dots+S_n}{t_1+t_2+\dots+t_n}.
\]

\textbf{Средняя скорость обычно не равна среднему арифметическому скоростей.}

\subsection*{Случай 1. Равные промежутки времени}
Если половину всего времени ехали со скоростью \(84\) км/ч, а вторую половину — \(56\) км/ч, то при общем времени \(T\):
\[
S_1=84\cdot\frac{T}{2}=42T,
\qquad
S_2=56\cdot\frac{T}{2}=28T.
\]
Поэтому
\[
v_{\text{ср}}=\frac{42T+28T}{T}=\boxed{70\text{ км/ч}}.
\]
Для двух равных интервалов времени действительно получается среднее арифметическое:
\[
v_{\text{ср}}=\frac{v_1+v_2}{2}.
\]

\subsection*{Случай 2. Равные участки пути}
Если половину трассы проехали со скоростью \(56\) км/ч, а вторую половину — \(84\) км/ч, то при длине трассы \(S\):
\[
t_1=\frac{S}{112},
\qquad
t_2=\frac{S}{168}.
\]
Следовательно,
\[
v_{\text{ср}}=
\frac{S}{\frac{S}{112}+\frac{S}{168}}
=\boxed{67{,}2\text{ км/ч}}.
\]
Для двух равных по длине участков можно использовать формулу:
\[
\boxed{v_{\text{ср}}=\frac{2v_1v_2}{v_1+v_2}}.
\]

\textbf{Главное различие:}
\begin{itemize}
  \item половина \textbf{времени} \(\Rightarrow\) равны времена, различаются пути;
  \item половина \textbf{трассы} \(\Rightarrow\) равны пути, различаются времена.
\end{itemize}

\section*{5. Многоэтажные дроби: быстрый технический приём}

При вычислениях средней скорости часто возникают сложные дроби. Полезно помнить:
\[
\frac{\frac{a}{b}}{\frac{c}{d}}=\frac{ad}{bc},
\qquad
\frac{\frac{a}{b}}{c}=\frac{a}{bc},
\qquad
\frac{a}{\frac{b}{c}}=\frac{ac}{b},
\]
где все знаменатели должны быть ненулевыми.

\textbf{Самопроверка:} деление на дробь всегда можно заменить умножением на обратную дробь.

\clearpage
\section*{Тренировочный блок — Путь А}
\textbf{10 задач.} Решения не записаны; после выполнения сверьтесь только с таблицей ответов.

\begin{enumerate}
  \item Два велосипедиста одновременно выехали из одной точки круговой трассы длиной \(18\) км в одном направлении. Их скорости равны \(72\) и \(54\) км/ч. Через сколько минут быстрый впервые обгонит медленного на один круг?

  \item Длина круговой трассы \(24\) км. Два участника стартуют одновременно из диаметрально противоположных точек и движутся в одном направлении. Быстрый участник находится сзади по ходу движения. Их скорости равны \(60\) и \(48\) км/ч. Через сколько минут быстрый догонит медленного?

  \item Катер проходит \(24\) км против течения на \(3\) часа дольше, чем те же \(24\) км по течению. Скорость течения \(2\) км/ч. Найдите собственную скорость катера.

  \item Теплоход идёт по течению со скоростью \(24\) км/ч, а против течения — \(16\) км/ч. На стоянку в пункте назначения ушло \(3\) часа, а весь рейс туда и обратно занял \(13\) часов. Найдите длину всего рейса.

  \item Скорость течения реки \(3\) км/ч. Плот вышел в путь на \(1\) час раньше моторной лодки и до общего момента окончания движения прошёл \(21\) км. Лодка за своё время прошла \(24\) км по течению и \(24\) км против течения. Найдите собственную скорость лодки.

  \item Первую половину всего времени автомобиль ехал со скоростью \(72\) км/ч, вторую половину — со скоростью \(48\) км/ч. Найдите среднюю скорость на всём пути.

  \item Первую половину трассы автомобиль проехал со скоростью \(60\) км/ч, вторую половину — со скоростью \(40\) км/ч. Найдите среднюю скорость.

  \item Турист прошёл \(10\) км со скоростью \(5\) км/ч, затем \(12\) км со скоростью \(6\) км/ч и ещё \(8\) км со скоростью \(4\) км/ч. Найдите среднюю скорость на всём маршруте.

  \item Половину трассы автомобиль ехал со скоростью \(72\) км/ч, вторую половину — со скоростью \(48\) км/ч. Найдите среднюю скорость. Ответ запишите в км/ч.

  \item Собственная скорость лодки \(10\) км/ч, скорость течения \(2\) км/ч. Лодка прошла \(24\) км по течению, сделала остановку на \(1\) час и вернулась \(24\) км против течения. Найдите среднюю скорость за весь рейс, включая остановку.
\end{enumerate}

\clearpage
\section*{Ответы}
\renewcommand{\arraystretch}{1.35}
\begin{center}
\begin{tabularx}{\linewidth}{@{}cX@{}}
\toprule
\textbf{№} & \textbf{Ответ} \\
\midrule
1 & \(60\) мин \\
2 & \(60\) мин \\
3 & \(6\) км/ч \\
4 & \(192\) км \\
5 & \(9\) км/ч \\
6 & \(60\) км/ч \\
7 & \(48\) км/ч \\
8 & \(5\) км/ч \\
9 & \(57{,}6\) км/ч \\
10 & \(8\) км/ч \\
\bottomrule
\end{tabularx}
\end{center}

\vfill
\begin{center}
{\small Лёвин Артём Александрович эксклюзивно для Анастасии Павловой}
\end{center}

\end{document}
